\section{Related Work}

\textbf{Evaluation validity under nuisance factors.} A line of work shows that reported LLM
capability is sensitive to factors that ought to be irrelevant. \citet{sclar2024formatspread}~(\emph{FormatSpread})
quantify how prompt-formatting choices alone move benchmark scores by margins comparable to
model differences; \citet{mizrahi2024stateofwhatart}~(\emph{State of What Art?}) show single-prompt evaluation is
unreliable and argue for multi-prompt protocols; \citet{dehghani2021benchmarklottery}~(\emph{The Benchmark Lottery})
document how benchmark and configuration choices determine which method appears to win.
Our contribution is adjacent but distinct: those works vary an input the evaluator controls
and observe score movement. \textbf{We identify a component of the measurement apparatus --- the
serving stack's tool-call parser --- that can zero out a capability signal entirely, without
any error, and whose distortion grows with model scale.}

\textbf{Constrained decoding and its costs.} \citet{tam2024speakfreely}~(\emph{Let Me Speak Freely?}) report that
format restrictions and structured-output constraints can degrade reasoning. This bears
directly on \S5.3, where enabling \texttt{strict:\ true} eliminates a 23\% role-confusion failure and
\emph{improves} pass rate (49 $\to$ 61). We do not claim constraints are free: our result is that on
this task the constraint's benefit --- admitting 23 previously-discarded items into execution
--- outweighs any reasoning cost, and we report the CoT control (\S5.3) where added reasoning
made matters worse. Whether constraints hurt elsewhere in the same stack is untested here.
Guided-decoding machinery \citep{willard2023outlines} underlies both \texttt{strict} and vLLM's
\texttt{required} mode.

\textbf{Tool-use benchmarks.} BFCL~\citep{bfcl2025}, $\tau$-bench \citep{yao2025taubench}, ToolLLM~\citep{qin2024toolllm},
API-Bank~\citep{li2023apibank} and Toolformer~\citep{schick2023toolformer} evaluate or train tool use. To our
knowledge these report tool-call rates as parsed by the serving layer. \S5.2 shows that
number can be zero while 80/100 well-formed calls are emitted, which --- if the pattern
replicates on their stacks --- would affect any absolute tool-use rate they report.

\textbf{Protocols and agent loops.} ReAct~\citep{yao2023react} is the text protocol we compare against.
vLLM~\citep{kwon2023vllm} is the serving stack; its tool-calling documentation specifies distinct
per-family configuration, which \S5.1 shows is not optional. verl~\citep{sheng2025hybridflow} provides the RL training stack; importantly its AgentLoop performs \textbf{its own}
tool-call parsing (\texttt{verl/experimental/agent\_loop/tool\_parser.py}, adapted from vLLM v0.9.1),
so a vLLM-side parser plugin does not affect training rollouts --- a distinction that cost us
one wasted experiment design (\S5.6).

\textbf{Prior report of the Qwen case.} vLLM issue \#32926~\citep{vllm_issue_32926},
opened 2026-01-23, documents that Qwen2.5-Coder was not trained on tool-call tokens, that it
emits \texttt{json} fenced blocks or bare JSON absent format instructions, and that a
\texttt{<tools>} few-shot template plus a dedicated parser reaches 98--100\% compliance. The
issue was closed as \emph{not planned} and labelled \emph{stale}. Section~\ref{sec:scale}
reproduces this baseline exactly---zero \texttt{<tool\_call>} and \texttt{<tools>} tags at
every scale---and extends it to scale quantification, to three further families, and to
training. A related verl issue~\citep{verl_issue_4124} reports the same role-confusion
signature we characterise in Section~\ref{sec:llama}: the model calls the task function
(\texttt{solve\_equation}) and the agent loop raises \texttt{KeyError}.

\textbf{RL with tool feedback.} GRPO~\citep{shao2024deepseekmath} and RLOO~\citep{ahmadian2024rloo} are the estimators
we train with; EvalPlus~\citep{liu2023evalplus} and KodCode~\citep{kodcode2025} supply held-out evaluation and training
problems respectively. We are not aware of prior work measuring interface-layer censoring
of tool trajectories inside an RL rollout, nor connecting it to which capability the policy
gradient can reach.
